%!TeX program = xelatex
% 沿用 SYSUReport 模板，作者 NorthSecond，CC BY 4.0。
\makeatletter
\def\input@path{{../../report-stage2/}}
\makeatother
\documentclass{SYSUReport}
\usepackage{array,tabularx,xurl}
\graphicspath{{../../report-stage2/figures/}}
\setmainfont{Times New Roman}
\setCJKmainfont[AutoFakeBold=2,AutoFakeSlant=0.15]{SimSun}
\setCJKsansfont{SimHei}
\setCJKmonofont{SimSun}
\setmonofont{Consolas}
\setlength{\headheight}{16pt}
\setlength{\parskip}{0.22em}
\renewcommand{\arraystretch}{1.2}
\newcolumntype{Y}{>{\raggedright\arraybackslash}X}
\newcommand{\code}[1]{\expandafter\nolinkurl\expandafter{\detokenize{#1}}}

\headl{数据库原理实验}
\headc{}
\headr{第10--13周阶段报告}
\lessonTitle{数据库原理实验报告}
\reportTitle{二次元周边店数据库编程与事务处理}
\stuname{郑立峰\quad 李维政}
\stuid{24325341\quad 24325143}
\inst{计算机学院}
\major{计算机科学与技术}
\date{2026年11月26日}
\hypersetup{pdftitle={二次元周边店数据库编程与事务处理},pdfauthor={郑立峰 李维政},hypertexnames=false}

\begin{document}
\hypersetup{pageanchor=false}
\cover\thispagestyle{empty}\clearpage
\pagenumbering{arabic}\hypersetup{pageanchor=true}

\section{实验目的和分工}

前一阶段的网页已经可以调整库存和完成现货结账。这一阶段把库存调整、销售统计、订单完成和补货入库写成可独立调用的数据库对象。调整数量时直接使用存储过程，库存触发器负责记录变化，查询则用函数、窗口查询和递归查询回答不同的问题。最后用两个数据库会话观察一笔业务尚未提交时，另一笔业务会看到什么。

这次 A 负责销售与订单，包括成交统计、结账触发器、窗口查询、订单事务和第13周双会话实验。B 负责库存调整、补货量、审计触发器、补货入库、商品分类递归和阶段整理。当前草稿记录 B 的实现与可取得的结果，销售和双会话章节随 A 的实测证据补齐。

\section{库存存储过程和函数}

店员登记订单时不会改变库存，库存员调整数量时则需要立即留下一条可核对的流水。网页把商品、增减量、原因和警戒值交给 \code{AdjustInventory}，员工编号来自登录身份。存储过程找到商品和库存记录，在同一事务中修改在手量和警戒值。调整后数量要能覆盖已有预留，调整量为零、商品不存在或员工无效时，过程返回对应错误。只需修改警戒值时，调用 \code{SetInventoryReorderPoint}，因此表单保留了单独修改警戒值的入口。

补货建议按可售量计算。例如 P 的在手量为 8，其中 3 件已经预留，实际可卖 5 件。若希望补到 12 件，建议量为 7。函数支持传入目标量，省略目标时用警戒值。达到目标返回 0，商品编号查不到时返回 \code{NULL}。库存页直接显示数据库函数的计算值。

\[
  \text{建议补货量}=\max\{\,\text{目标量}-(\text{在手量}-\text{预留量}),\,0\,\}.
\]

\section{库存流水触发器}

一次补货或扣库可能同时涉及 P、Q 两种商品。所以触发器将 \code{inserted} 和 \code{deleted} 按商品号对应起来，用一条查询计算全部变动行的前后值。调整警戒值时两种数量都相等，库存表会保存新警戒值，流水表的行数保持原值。插入与删除按缺少的一侧数量为 0 计算差额。

同一个数据库连接在写库存前提供业务类型、来源单号和操作人。手工调整另外提供原因，序列生成的调整编号作为来源。销售与补货可按单号查找对应的商品流水。触发器与扣库或加库在同一事务内执行，所以回滚业务时，该次流水也一起回滚。

\begin{center}
\begin{tabularx}{\textwidth}{lYY}
\toprule
类型 & 流水的来源 & 操作人 \\
\midrule
销售 & 销售订单的内部编号 & 实际执行结账的员工 \\
补货 & 采购单的内部编号 & 执行验收的员工 \\
手工调整 & 调整序列号，并记录原因 & 当前登录的员工 \\
直接变更 & SQL 维护操作 & 数据库的原始登录名 \\
\bottomrule
\end{tabularx}
\end{center}

\section{补货事务和分类递归}

\subsection{整单入库}

补货沿用原采购单表和明细表。人工先登记本次订购的商品和数量。验收时输入订单号与员工编号，\code{ReceiveRestock} 先锁住该采购单，检查类型为 \code{RESTOCK}、状态为 \code{ORDERED}，并核对各明细尚未验收。接着按商品号从小到大锁住库存并加库。每项都成功后，明细的实收量设为订购量，表头记录 \code{RECEIVED}、验收时间和经办人。

同一张单再次入库会先取得同一表头的锁，看到 \code{RECEIVED} 后直接结束。测试中安排编号较大的 Q 在入库时超出整数库存上限。这样第一次加库与对应流水已在事务内写入，第二项才报错，可以核对先前写入的 P 是否被撤回，并检查单据、明细与流水是否同时恢复。

\subsection{分类子树}

分类表增加一个可空的父级编号。查询以指定分类为第0层，每轮用上一层的分类号连接子分类的父级号，同时把子分类号接到路径上。以「周边」为起点，沿「展示类」走到「亚克力立牌」，深度为2。另一条分支从「随身类」走到「钥匙扣」。路径中的编号用于跳过重复节点，查询还设有100层递归上限。查询叶子时只返回起点，查询不存在的编号时结果为空。

\section{并发实验交接}

共用基线设 P 在手为 10、Q 为 8，两者预留都为 0。补货用单据 \code{90013101} 给 P 入库 5 件。通过外层事务调用真实的补货过程，可以让库存、流水、明细和单据一起停在提交前，用另一个会话观察读取或等待。若订单先扣 P 的 2 件并提交，随后补货的预期库存为 13。若订单回滚，补货后的预期库存为 15。

两种业务都按同一商品号顺序加锁。第13周脚本需要记录会话号、锁资源、等待状态与最终数据，分别比较提交、回滚和反向访问时的现象。死锁和防超卖结论在订单过程及两会话实验结果加入后填写。

\section{验证记录与阶段小结}

本次运行了18项 Web 调用测试和6项密码辅助回归，全部通过。Web 测试通过 Flask 的真实路由发起请求，并用驱动替身读取调用参数。举例说，请求额外传入员工编号99时，库存过程仍收到登录绑定的员工编号2。测试也核对了库存不足、非法数量、跨角色调用以及 v1.0 与第三阶段配置下的不同入口。

11个本阶段 SQL 文件通过 SQL Server ScriptDom 语法解析。11月27日，在 SQL Server 17.0.1000.7 的独立验收库执行三组断言，库存、补货及分类递归均返回 \code{PASS}。库存组以对应的库内测试角色调用调整过程和补货量函数。

同日用库存员账户在 Chrome 中核对了 P003 的补货建议。初始库存为2件，预留为0，警戒值为3，页面建议补1件。盘盈1件后，库存达到警戒值，建议量降为0。单独把警戒值提高到5，库存仍为3，页面建议补2件，刷新后读数一致。最后将库存和警戒值恢复为2和3。数据库中该商品共有两条调整流水，对应盘盈和恢复操作，均记录了员工编号2及表单填写的原因。单独修改警戒值没有新增库存流水。

\begin{center}
\begin{tabularx}{\textwidth}{p{3.4cm}YY}
\toprule
操作 & 初始值和输入 & 实测结果 \\
\midrule
调整 P & 在手10，调整$-2$ & 在手8，一条调减2的流水 \\
批量扣库回滚 & P、Q 各在手8，分别扣2 & 事务内两条同来源流水，回滚后库存合计16且两条流水撤回 \\
验收入库并重复调用 & P=10，补货单为 P$\times$5 & 第一次 P=15，明细实收5，表头已验收。重复调用报53112 \\
整单中途失败 & P=15，Q=2147483646，订购 P$\times$5、Q$\times$3 & 报53105，P=15、Q 保持原值，单据待验收，明细实收0，本单流水0条。恢复 Q 后重试得到 P=20、Q=11 \\
分类子树 & 根、中间、叶和不存在编号 & 返回5、2、1、0行 \\
\bottomrule
\end{tabularx}
\end{center}

复验时，单独读取 \code{XACT_STATE()} 得到0，在含表查询的语句中读取时得到1。在表查询之前单独读取事务状态，重跑断言通过。外层事务暂停点也经过单会话验证。入库后事务内 P=15，单据已验收并写入一条流水。会话回滚后 P=10，单据恢复待验收，本单流水为0条。

\end{document}
